% How Many Sensors Are Enough? An Exhaustive Sensor-Subset Analysis for
% Rotating Machinery Fault Diagnosis
%
% IEEEtran conference skeleton.
%
% IEEEtran.cls is NOT included in this repository. If your TeX installation does
% not provide it, get it from CTAN (https://ctan.org/pkg/ieeetran) or compile on
% Overleaf, which ships it.
%
% Build:
%   pdflatex main.tex && bibtex main && pdflatex main.tex && pdflatex main.tex
%
% EVERY RESULT SENTENCE BELOW IS A PLACEHOLDER. Fill them in from
% ../tables/hypotheses.txt and ../tables/table_main.csv. Do not type a number
% that you have not copied from a generated file.

\documentclass[conference]{IEEEtran}

\usepackage{booktabs}
\usepackage{graphicx}
\usepackage{amsmath}
\usepackage{cite}
\usepackage{url}          % the dataset URL contains a tilde

\graphicspath{{../figures/}}

% Figure captions are GENERATED by 04_analyze.py so that what they claim about
% aggregation always matches what the code did. Regenerate, do not hand-edit.
\input{../figures/captions.tex}

\begin{document}

\title{How Many Sensors Are Enough? An Exhaustive Sensor-Subset Analysis for
Rotating Machinery Fault Diagnosis}

\author{\IEEEauthorblockN{TODO Author Name}
\IEEEauthorblockA{TODO Department\\
TODO Institution\\
TODO email}}

\maketitle

\begin{abstract}
% TODO: one paragraph. State the question (how does diagnosis performance depend
% on the number and choice of sensors), the method (one fixed Random Forest over
% all 127 subsets of 7 sensors, leakage-free grouped cross-validation, selection
% on training data only), and the finding. No numbers until they come from
% ../tables/.
TODO abstract.
\end{abstract}

\begin{IEEEkeywords}
fault diagnosis, rotating machinery, sensor selection, feature importance,
machine learning, MaFaulDa
\end{IEEEkeywords}

\section{Introduction}
% TODO: why sensor count matters in practice (cost, wiring, installation,
% maintenance, retrofit).
% TODO: the gap. Most work adds sensors or proposes methods; few measure how
% performance degrades as sensors are removed, with selection done honestly.
% TODO: state RQ1-RQ3 and H1-H3.
% TODO: contributions. Exhaustive rather than heuristic subset evaluation;
% leakage-free grouped CV; selection on training data only; pre-registered
% hypotheses; fully reproducible pipeline.

\section{Related Work}
% Three subsections, fed by the five reading bins in refs.bib. Keep the
% bin-to-subsection mapping in step with the header of that file:
%   bins 1 + 2 -> II-A,  bin 3 -> II-B,  bins 4 + 5 -> II-C.
% Budget roughly 0.75 of a column in total, so about one paragraph each plus
% the closing gap paragraph.

\subsection{Fault Diagnosis on Rotating Machinery}
% Fed by refs.bib bins 1 and 2.
% TODO: what is conventionally measured, which feature families are standard,
% and what performance is typically reported.
% TODO: prior work on the MaFaulDa dataset \cite{mafaulda} specifically. For
% each reported number, state the validation scheme it was obtained under; a
% number from a window-level random split is not comparable with ours.

\subsection{Sensor and Feature Selection and Reduction}
% Fed by refs.bib bin 3.
% TODO: wrapper, filter and embedded selection; greedy forward selection in
% particular.
% TODO: whether any prior work evaluates sensor subsets exhaustively rather
% than heuristically, which is what separates this study's design.

\subsection{Evaluation Methodology in Machinery Datasets}
% Fed by refs.bib bins 4 and 5. This subsection carries real weight: it is
% where this paper's leakage handling and its treatment of permutation
% importance are positioned against prior work. It is not filler.
% TODO: window-level and segment-level leakage, grouped cross-validation, and
% reported cases of inflated results from splitting correlated segments at
% random. Position the leakage check in Section III against these.
% TODO: permutation importance and its known behaviour under correlated or
% redundant features: it measures MARGINAL contribution, so a near-zero score
% cannot on its own separate a sensor that carries no information from one
% that is merely redundant. Position our use of forward selection and the
% sensor correlation matrices as the response to that limitation.

% TODO: closing paragraph. State the gap this paper addresses: prior work
% largely adds sensors or proposes methods; few measure how performance
% degrades as sensors are removed, with selection performed on training data
% only and with the evaluation protected against leakage.

\section{Dataset and Methodology}

\subsection{Dataset}
% TODO: MaFaulDa description \cite{mafaulda}: 1951 recordings, 6 classes,
% 8 channels recorded, 7 used as sensors, 50 kHz, 5 s per recording, speed
% range. State the tachometer exclusion here and cross-reference Section V.
% TODO: state that the class label is the top-level folder and that subtype and
% severity are recorded but not used as labels.
%
% TODO: DEFINE THE TWO BEARING CLASSES ONCE, HERE, FOR THE READER.
% "Underhang" and "overhang" denote bearing POSITIONS relative to the rotor,
% not defect types. Defect types (ball, cage, outer race) appear in the
% subfolders beneath them and are merged into the two position classes by this
% study's 6-class choice. Note also that the dataset's own web page words these
% same two classes differently, so a reader coming from that page needs the
% mapping stated explicitly. See PROTOCOL.md section 2.1.

\subsection{Feature Extraction}
% TODO: 5 windows per recording, 18 features per sensor per window, identical
% for every sensor so that a subset's score reflects the sensors and not a
% differing feature count.
% TODO: state explicitly that standard deviation is excluded because after DC
% removal it equals RMS.

\subsection{Cross-Validation and Leakage Control}
% TODO: StratifiedGroupKFold, 5 folds x 2 repeats, grouped by recording.
% TODO: report the leakage check from 02_checkpoint.py: the naive random-window
% split against the grouped split, and the inflation between them.
% TODO: state that test-fold isolation is enforced by assertions, not convention.

\subsection{Scoring}
% TODO: macro-F1 as the primary metric.
% TODO: the per-fold label set: one label set per outer fold, taken from y_true,
% shared by every subset in that fold. Explain why the scikit-learn default is
% unsuitable here (the denominator would vary between subsets within a fold).

\subsection{Sensor Selection}
% TODO: grouped permutation importance, one shared permutation across all 18
% columns of a sensor, common random numbers across sensors.
% TODO: greedy forward selection.
% TODO: state that both use the training fold only, via an inner 75/25 split.

\subsection{Exhaustive Evaluation}
% TODO: all 127 subsets on every split; 1270 score rows.
% TODO: state that the tachometer is excluded from the sensor set and why
% (label-correlated recording artefact); see PROTOCOL.md section 2.2. Note
% that its features are still extracted and that the 8-sensor comparison is
% reported separately.
% TODO: state that selected subsets are looked up rather than refit, and why
% that is exact.

\section{Results}

% TODO: RQ1. How macro-F1 changes as k falls from 8 to 1.
\input{../tables/table_main.tex}

\begin{figure}[t]
  \centering
  \includegraphics[width=\columnwidth]{fig1_f1_vs_k.pdf}
  \caption{\FigOneCaption}
  \label{fig:f1_vs_k}
\end{figure}

% TODO: H1 verdict, from ../tables/hypotheses.txt. Report the PRIMARY verdict
% (the fair curves) and the SECONDARY oracle verdict separately, and state that
% the oracle curve is a maximum over C(8,k) candidates whose upward bias peaks
% at k=4.

% TODO: H2 verdict. Report the FAIR k as the answer. The oracle k is an upper
% bound and must be labelled as such.

\begin{figure}[t]
  \centering
  \includegraphics[width=\columnwidth]{fig2_sensor_importance.pdf}
  \caption{\FigTwoCaption}
  \label{fig:importance}
\end{figure}

% TODO: RQ3 and H3. Which sensors rank highest, and whether informed selection
% beats the median subset of the same size.

\begin{figure}[t]
  \centering
  \includegraphics[width=\columnwidth]{fig3_per_class_recall.pdf}
  \caption{\FigThreeCaption}
  \label{fig:per_class}
\end{figure}

\begin{figure}[t]
  \centering
  \includegraphics[width=\columnwidth]{fig4_detection_errors.pdf}
  \caption{\FigFourCaption}
  \label{fig:errors}
\end{figure}

\begin{figure*}[t]
  \centering
  \includegraphics[width=\textwidth]{fig5_confusion.pdf}
  \caption{\FigFiveCaption}
  \label{fig:confusion}
\end{figure*}

\begin{figure}[t]
  \centering
  \includegraphics[width=\columnwidth]{fig6_sensor_correlation.pdf}
  \caption{\FigSixCaption}
  \label{fig:sensor_corr}
\end{figure}

% Figures 6 and 6b are alternative views of the same redundancy question. Keep
% whichever reads better on the real data; both are descriptive and enter no
% hypothesis test.
\begin{figure}[t]
  \centering
  \includegraphics[width=\columnwidth]{fig6b_mean_abs_correlation.pdf}
  \caption{\FigSixBCaption}
  \label{fig:sensor_corr_mean}
\end{figure}

\subsection{Secondary Analysis: The Excluded Tachometer Channel}
% TODO: run only AFTER the 7-sensor study is complete. Repeat the whole
% pipeline with the tachometer included (8 sensors, 255 subsets) and report
% where it ranks and what it does to the curve. If it ranks highly, that is a
% concrete demonstration of an artefact-driven shortcut: a channel that carries
% no fault information can still dominate a sensor ranking because its
% recording artefact is label-correlated. Present it as a contribution and a
% caution for other sensor-selection work, not as a result of this study.
% See PROTOCOL.md section 2.2.

\section{Discussion}
% TODO: what the sensor-count curve implies for instrumentation decisions.
% TODO: which fault classes degrade first as sensors are removed, and what that
% means for a practitioner choosing sensors for a specific failure mode.
% TODO: redundancy between channels, read from Fig. \ref{fig:sensor_corr} and
% Fig. \ref{fig:sensor_corr_mean} together with the forward-selection order.
% TODO: interpret low-importance sensors carefully. Permutation importance
% measures marginal contribution, so a low value may indicate redundancy rather
% than absence of information; say which interpretation the correlation matrix
% and the forward-selection order support.

\subsection{Limitations}
% TODO: expand each of the following into prose. Do not delete any of them.
%
% 1. Single dataset. All results come from MaFaulDa. Generalisation to other
%    machines, mountings and fault types is untested.
% 2. Laboratory rig. The recordings come from a controlled test rig, not from
%    machinery in service, so the noise, loading and mounting conditions are
%    more benign than in the field.
% 3. One classifier. A single fixed Random Forest is used throughout, by design.
%    Conclusions about sensor value may not transfer to other model families.
% 4. Design-time sensor reduction, not runtime sensor loss. Every subset model
%    is retrained on that subset. This answers "which sensors should we install",
%    not "what happens when an installed sensor fails mid-operation", which would
%    require evaluating a full-sensor model under missing inputs.
% 5. Non-independent scores. The 10 cross-validation scores come from 2 repeats
%    over the same recordings, so they are not mutually independent. The
%    p-values are supporting evidence alongside effect sizes and percentile
%    ranks, not standalone proof.
%
% 6. Tachometer channel quality on a subset of recordings. On 67 of the 1,951
%    recordings, all of them ball_fault, a rising-edge estimate of rotation
%    speed from the tachometer channel overcounts, reporting up to 320.6 Hz
%    against a true 50.79 Hz on a rig specified at 700-3600 rpm (about
%    12-60 Hz). Those traces swing about 22 percent wider (median standard
%    deviation 2.118 against 1.742), with a more negative minimum and a lower
%    maximum, which lowers the midpoint threshold from about 1.97 to 0.80-1.27
%    and into the noise band. An FFT of the same channel recovers the correct
%    frequency to within about 1.3 Hz, so the signal carries the information
%    and only the threshold crossing fails. This study therefore takes the
%    rotation frequency from the file name, which the dataset encodes, and
%    reports the tachometer estimate alongside it without using it. The
%    limitation is a property of the dataset's tachometer channel on those
%    files rather than of the method, and it matters for any work that derives
%    speed from that channel. See PROTOCOL.md section 3.2 and
%    scripts/diagnose_speed.py.

\section{Conclusion}
% TODO: answer RQ1-RQ3 in one paragraph, using only numbers from ../tables/.
% TODO: state the fair H2 k as the headline answer, never the oracle k.
% TODO: future work.

\section*{Acknowledgments}
% TODO: funding and people.
%
% AI ASSISTANCE DISCLOSURE (required by IEEE policy; do not delete).
% IEEE requires that the use of AI-generated text or code be disclosed in the
% acknowledgments. Adapt the following to what was actually used, and keep it
% accurate:
%
% "The author used [tool name and version] to assist with [the implementation of
% the analysis pipeline / drafting and editing portions of the manuscript]. All
% experimental design decisions, result interpretation and conclusions are the
% author's own, and the author reviewed and verified all generated content and
% takes full responsibility for it."
TODO acknowledgments, including the AI assistance disclosure above.

\bibliographystyle{IEEEtran}
\bibliography{refs}

\end{document}
